\chapter*{Tóm tắt báo cáo}
\addcontentsline{toc}{chapter}{Tóm tắt báo cáo}

Báo cáo trình bày thiết kế và hiện thực một nguyên mẫu gồm hai lớp kiểm soát
bổ sung cho nhau. Lớp phòng ngừa sử dụng Role-Based Access Control (RBAC) để
buộc quyền tại backend của ứng dụng Nova Bank và minh họa separation of duties
giữa Teller và Controller. Lớp giám sát chuẩn hóa nhật ký thành Event, áp dụng
luật giải thích được để phát hiện chỉ báo SQL Injection, password guessing và
truy cập trái phép, sau đó chấm điểm, tương quan và xuất bằng chứng dưới dạng
CSV/JSON. Audit log của Nova Bank được nối với pipeline bằng một adapter thay
vì chỉ được mô tả như hướng phát triển.

Đánh giá gồm ba tầng. Bộ regression 60 Event kiểm tra tính ổn định của các
scenario ban đầu. Bộ holdout tổng hợp 300 Event được sinh bằng quy trình và
manifest riêng, gồm 180 benign, 40 SQL Injection, 40 password guessing và 40
unauthorized access; trong đó 60 Event là hard negative. Trên holdout, macro
precision đạt 1,0000, macro recall 0,9333 và macro F1 0,9630. SQL Injection
bỏ sót 8 biến thể encoding/obfuscation; password guessing phát hiện đủ 8 chiến
dịch, nhưng phân tích độ nhạy cho thấy kết quả thay đổi mạnh khi ngưỡng chuyển
từ 3 sang 5 hoặc 7 lần thất bại. Bộ context-risk 14 Event được dùng để kiểm
chứng khả năng giữ quan hệ Event--Alert--Incident, chưa dùng để tuyên bố độ
chính xác thống kê.

Kết quả chứng minh nguyên mẫu có thể tái lập, có enforcement tại backend và có
đường truy vết từ quyết định nghiệp vụ tới cảnh báo. Tuy nhiên, toàn bộ dữ liệu
đánh giá vẫn là dữ liệu tổng hợp, holdout chỉ độc lập về quy trình chứ chưa độc
lập về nguồn vận hành, và risk score là thang ưu tiên theo luật chứ không phải
xác suất tấn công. Vì vậy kết luận phù hợp nhất là hệ thống đáp ứng mục tiêu
đồ án và tạo nền tảng thử nghiệm có kiểm soát; trước khi triển khai thực tế cần
đánh giá trên log đã ẩn danh, có nhãn độc lập, nhiều thời kỳ và baseline cạnh
tranh.
